%%
%% design.tex
%% 
%% Made by Alex Nelson <pqnelson@gmail.com>
%% Login   <alex@lisp>
%% 
%% Started on  2026-06-02T09:21:47-0700
%% Last update 2026-06-02T09:21:47-0700
%% 

\section{Design of the proof assistant}

\begin{node}
The design of a proof assistant requires a few design
choices. LCF-style proof assistants share a similar design, which
simplifies our discussion: what is the design of Myo anyways?
\end{node}

\begin{aside}[What is ``software design'' anyways?]
Discussions of ``software design'' tend to be littered with strong
opinions, even with what qualifies as ``design''. We should
distinguish between the activity of ``software design'' as a component
of programming, from the ``design'' of a particular piece of
software. We will not even attempt to answer the question, ``What is
software design in the context of programming?'' Instead we focus on
the question, ``What is the design of an LCF-style proof assistant?''

We will broadly discuss the \textbf{architecture} without worrying too
much about the minutiae. The reader can think of this section as
discussing the milestones expected for any proof assistant.
\end{aside}

\begin{node}[Syntax trees]
The first thing we need to do is to represent terms and formulas using
data structures.

We do this by implementing a \texttt{Term} module with a
\texttt{Term.t} type, which is just an algebraic data type.

Similarly, we have a \texttt{Formula} module with a \texttt{Formula.t}
abstract data type. The code for constructing, destructing, and
determining the type of formula is contained in this module.

This step is standard for all proof assistants, interpreters,
compilers, and so on: we need to get a computer representation of the
input, some sort of abstract syntax tree. There is a lot I can say
about this (literally, I can ramble on for pages), but suffices to say
it is the first step on any proof assistant.
\end{node}

\begin{node}[Prelogic]
At this step, we have the data structures implemented for the syntax
trees of terms and formulas. We can follow Slind~\cite{slind1990implementation}
discussion of implementing \HOL\ and implement the ``prelogic'' for
our first-order logic. \HOL\ needs to worry about $\alpha$-congruence,
term substitution, and type substitution.

Fortunately, first-order logic only needs to address substituting a
term for a variable in either a formula or a term.
\end{node}

\begin{node}[Kernel]
Once we have data structures for terms and formulas, the next thing we
will want to do is implement the ``kernel''. For LCF-style proof
assistants based on (propositional, first-order, or higher-order)
logic, this means implementing a data type for theorems which is an
opaque type (or, for object-oriented languages, have private
constructors). The only ways to construct theorems would be using
axioms (functions taking terms or formulas and produce theorems) and
inference rules (functions which take one or more theorems and produce
a new theorem).

At this point, the only possible threats to soundness will be in the
code for the syntax trees or the kernel. This is why LCF-style proof
assistants want as small a kernel as possible.
\end{node}

\begin{node}[Derived inference rules and theorems]
Although we have obtained something which qualifies as a ``proof assistant'',
there is still much left to be desired.
What we want is a \emph{usable} proof assistant.

The first thing we need is a library of theorems and derived inference
rules. This may be the least exciting part of the proof assistant,
because we just enumerate things in the style of Homer's catalog of ships.
\end{node}

\begin{node}[Tactics and procedural-style proofs]
We will eventually want declarative-style proofs like you'd find in
Mizar. The first step towards that end is to implement
procedural-style proofs using tactics and goals. ``Goals'' are the
proof obligations yet to be proven, ``tactics'' give us a way to
construct backwards-style proofs from the theorem statement to axioms
or assumptions. The disadvantage to procedural-style proofs is the
unreadability, and since proofs are read far more often than written,
this is a major disadvantage.

Declarative-style proofs are designed with readability in mind. We
will try to implement something approximating Mizar's proof steps, but
for Intuitionistic logic.
\end{node}

\begin{node}[Next steps]
What happens next? There is no shortage of possible avenues to
explore. The low hanging fruit will be to provide some sort of proof
procedure like tableaux or resolution, so ``obvious'' inferences can
be handled ``automatically''.

After that, definitions will be critically important. In first-order
logic, defining a new predicate $P[x_{1},\dots,x_{n}]$ amounts to
introducing the axiom $\vdash\forall\vec{x}\ldotp P[\vec{x}]\iff\Phi[\vec{x}]$
where $\Phi[\vec{x}]$ is some formula whose free variables form a
subset of $\{x_{1},\dots,x_{n}\}$. For terms, we either have
definitions of the form ``The term $t(\vec{x})$ means $P[\vec{x},t(\vec{x})]$ holds''
or we have an abbreviation. In the former case, we need to prove
the theorem $\vdash\forall\vec{x}\ldotp\exists!f\ldotp P[\vec{x},f(\vec{x})]$.
We also need a dictionary to store these definitions.

Term rewriting and simplification are desirable components for any
LCF-style proof assistant. It quickly becomes tedious manually proving
what term rewriting can do for us. Paulson~\cite{paulson1983higher} introduced
``conversions'' for this purpose.
\end{node}